% GENERATED FILE. Edit canonical lessons, then run npm run generate:books.
% canonical-source-hash: fnv1a64:b45e9d32905fd2d7
% canonical-lessons: AR-C293-murasil, AR-C293-bayan, AR-C293-sahifa, AR-C293-alana, AR-C293-sarraha

\chapter{Reporters and statements}
\label{ch:ar-reporters-and-statements}
\hlchaptermodality{293}

\begin{chapteropening}
\textbf{By the end of this chapter:} \emph{I can say a correspondent, a statement, a newspaper, to announce and to declare.}

Complete the last lesson of chapter 293: I can say a correspondent, a statement, a newspaper, to announce and to declare.
\end{chapteropening}

\section[\texorpdfstring{\ar{مراسل} (murāsil)}{مراسل (murāsil)}]{\ar{مراسل} (murāsil) — a correspondent, a reporter}

\label{lesson:AR-C293-murasil}



\begin{quote}
Before the new one: say the Arabic for a sentence, then the Arabic for a barber.
\end{quote}

\subsection*{You'll want to know: \ar{مراسل}}
\textbf{\ar{مراسل}} — \emph{murāsil} — \textquotedblleft{}a correspondent, a reporter\textquotedblright{}.

\begin{quote}
\ar{كتب} \ar{المراسل} \ar{عن} \ar{المطر} \ar{في} \ar{القرية}.
\end{quote}

\begin{quote}
\emph{kataba al-murāsil ʿan al-maṭar fī al-qarya.}
\end{quote}

\begin{quote}
The correspondent wrote about the rain in the village.
\end{quote}

\begin{cousinweb}
\textbf{\ar{مراسل}} (\emph{murāsil}) is built on the native Arabic root \emph{r-s-l} (Lane), the root of \emph{arsala}, \textquotedblleft{}to send\textquotedblright{}. \emph{Murāsil} is the Form III participle, \textquotedblleft{}one who writes back and forth\textquotedblright{}: a correspondent.
\end{cousinweb}

\begin{grammarlens}[title={What you've built}]
One more word for telling what happened.
\end{grammarlens}

\subsection*{Your turn}
\begin{itemize}
\raggedright
  \item \emph{Say it:} \emph{murāsil}
  \item \emph{Say it:} \emph{murāsil}, once more
  \item \emph{Say it:} \emph{kataba al-murāsil ʿan al-maṭar fī al-qarya.} — The correspondent wrote about the rain in the village.
  \item \emph{Recall:} say the Arabic for a sentence, then the Arabic for a barber, then say \emph{murāsil} again
\end{itemize}

\begin{tcolorbox}[breakable,colback=teal!4,colframe=teal!35!black,title={Before you move on}]
What does \ar{مراسل} mean? (\textquotedblleft{}A correspondent, a reporter\textquotedblright{}.) Say it once more.
\end{tcolorbox}
\section[\texorpdfstring{\ar{بيان} (bayān)}{بيان (bayān)}]{\ar{بيان} (bayān) — a statement; a clear account}

\label{lesson:AR-C293-bayan}



\begin{quote}
Before the new one: say the Arabic for a barber, then the Arabic for a correspondent.
\end{quote}

\subsection*{You'll want to know: \ar{بيان}}
\textbf{\ar{بيان}} — \emph{bayān} — \textquotedblleft{}a statement; a clear account\textquotedblright{}.

\begin{quote}
\ar{قرأ} \ar{المدير} \ar{بيانا} \ar{قصيرا}.
\end{quote}

\begin{quote}
\emph{qaraʾa al-mudīr bayānan qaṣīran.}
\end{quote}

\begin{quote}
The manager read out a short statement.
\end{quote}

\begin{cousinweb}
\textbf{\ar{بيان}} (\emph{bayān}) is built on the native Arabic root \emph{b-y-n} (Lane), \textquotedblleft{}to be apart; to become clear\textquotedblright{}, the root of \emph{bayna}, \textquotedblleft{}between\textquotedblright{}. A \emph{bayān} makes a thing clear.
\end{cousinweb}

\begin{grammarlens}[title={What you've built}]
One more word for telling what happened.
\end{grammarlens}

\subsection*{Your turn}
\begin{itemize}
\raggedright
  \item \emph{Say it:} \emph{bayān}
  \item \emph{Say it:} \emph{bayān}, once more
  \item \emph{Say it:} \emph{qaraʾa al-mudīr bayānan qaṣīran.} — The manager read out a short statement.
  \item \emph{Recall:} say the Arabic for a barber, then the Arabic for a correspondent, then say \emph{bayān} again
\end{itemize}

\begin{tcolorbox}[breakable,colback=teal!4,colframe=teal!35!black,title={Before you move on}]
What does \ar{بيان} mean? (\textquotedblleft{}A statement; a clear account\textquotedblright{}.) Say it once more.
\end{tcolorbox}
\section[\texorpdfstring{\ar{صحيفة} (ṣaḥīfa)}{صحيفة (ṣaḥīfa)}]{\ar{صحيفة} (ṣaḥīfa) — a gazette, a newspaper}

\label{lesson:AR-C293-sahifa}



\begin{quote}
Before the new one: say the Arabic for a correspondent, then the Arabic for a statement.
\end{quote}

\subsection*{You'll want to know: \ar{صحيفة}}
\textbf{\ar{صحيفة}} — \emph{ṣaḥīfa} — \textquotedblleft{}a gazette, a newspaper\textquotedblright{}.

\begin{quote}
\ar{تقرأ} \ar{أمي} \ar{صحيفة} \ar{كل} \ar{صباح}.
\end{quote}

\begin{quote}
\emph{taqraʾu ummī ṣaḥīfa kulla ṣabāḥ.}
\end{quote}

\begin{quote}
My mother reads a newspaper every morning.
\end{quote}

Plural: \ar{صحف}, \emph{ṣuḥuf}.

\begin{cousinweb}
\textbf{\ar{صحيفة}} (\emph{ṣaḥīfa}) is built on the native Arabic root \emph{ṣ-ḥ-f} (Lane), where it names a written leaf or page. Modern Arabic gave the old word the newspaper.
\end{cousinweb}

\begin{grammarlens}[title={What you've built}]
One more word for telling what happened.
\end{grammarlens}

\subsection*{Your turn}
\begin{itemize}
\raggedright
  \item \emph{Say it:} \emph{ṣaḥīfa}
  \item \emph{Say it:} \emph{ṣaḥīfa}, once more
  \item \emph{Say it:} \emph{taqraʾu ummī ṣaḥīfa kulla ṣabāḥ.} — My mother reads a newspaper every morning.
  \item \emph{Recall:} say the Arabic for a correspondent, then the Arabic for a statement, then say \emph{ṣaḥīfa} again
\end{itemize}

\begin{tcolorbox}[breakable,colback=teal!4,colframe=teal!35!black,title={Before you move on}]
What does \ar{صحيفة} mean? (\textquotedblleft{}A gazette, a newspaper\textquotedblright{}.) Say it once more.
\end{tcolorbox}
\section[\texorpdfstring{\ar{أعلن} (aʿlana)}{أعلن (aʿlana)}]{\ar{أعلن} (aʿlana) — to announce}

\label{lesson:AR-C293-alana}



\begin{quote}
Before the new one: say the Arabic for a statement, then the Arabic for a gazette.
\end{quote}

\subsection*{You'll want to know: \ar{أعلن}}
\textbf{\ar{أعلن}} — \emph{aʿlana} — \textquotedblleft{}to announce\textquotedblright{}.

\begin{quote}
\ar{أعلنت} \ar{المدرسة} \ar{موعد} \ar{الامتحان}.
\end{quote}

\begin{quote}
\emph{aʿlanat al-madrasa mawʿid al-imtiḥān.}
\end{quote}

\begin{quote}
The school announced the date of the exam.
\end{quote}

Present: \ar{يعلن}, \emph{yuʿlin}.

\begin{cousinweb}
\textbf{\ar{أعلن}} (\emph{aʿlana}) is built on the native Arabic root \emph{ʿ-l-n} (Lane), \textquotedblleft{}to be open, public\textquotedblright{}. It is Form IV, \textquotedblleft{}to make public\textquotedblright{}.
\end{cousinweb}

\begin{grammarlens}[title={What you've built}]
One more word for telling what happened.
\end{grammarlens}

\subsection*{Your turn}
\begin{itemize}
\raggedright
  \item \emph{Say it:} \emph{aʿlana}
  \item \emph{Say it:} \emph{aʿlana}, once more
  \item \emph{Say it:} \emph{aʿlanat al-madrasa mawʿid al-imtiḥān.} — The school announced the date of the exam.
  \item \emph{Recall:} say the Arabic for a statement, then the Arabic for a gazette, then say \emph{aʿlana} again
\end{itemize}

\begin{tcolorbox}[breakable,colback=teal!4,colframe=teal!35!black,title={Before you move on}]
What does \ar{أعلن} mean? (\textquotedblleft{}To announce\textquotedblright{}.) Say it once more.
\end{tcolorbox}
\section[\texorpdfstring{\ar{صرح} (ṣarraḥa)}{صرح (ṣarraḥa)}]{\ar{صرح} (ṣarraḥa) — to declare, to state}

\label{lesson:AR-C293-sarraha}



\begin{quote}
Before the new one: say the Arabic for a gazette, then the Arabic for to announce.
\end{quote}

\subsection*{You'll want to know: \ar{صرح}}
\textbf{\ar{صرح}} — \emph{ṣarraḥa} — \textquotedblleft{}to declare, to state\textquotedblright{}.

\begin{quote}
\ar{صرح} \ar{المدير} \ar{بأن} \ar{المكتب} \ar{مغلق}.
\end{quote}

\begin{quote}
\emph{ṣarraḥa al-mudīr bi-anna al-maktab mughlaq.}
\end{quote}

\begin{quote}
The manager stated that the office was closed.
\end{quote}

Present: \ar{يصرح}, \emph{yuṣarriḥ}.

It takes \ar{بـ} before what is stated: \ar{صرح} \ar{بأن} …

\begin{cousinweb}
\textbf{\ar{صرح}} (\emph{ṣarraḥa}) is built on the native Arabic root \emph{ṣ-r-ḥ} (Lane), \textquotedblleft{}pure, unmixed, clear\textquotedblright{}. It is Form II, \textquotedblleft{}to make clear\textquotedblright{}: to say a thing plainly.
\end{cousinweb}

\begin{grammarlens}[title={What you've built}]
One more word for telling what happened.
\end{grammarlens}

\subsection*{Your turn}
\begin{itemize}
\raggedright
  \item \emph{Say it:} \emph{ṣarraḥa}
  \item \emph{Say it:} \emph{ṣarraḥa}, once more
  \item \emph{Say it:} \emph{ṣarraḥa al-mudīr bi-anna al-maktab mughlaq.} — The manager stated that the office was closed.
  \item \emph{Recall:} say the Arabic for a gazette, then the Arabic for to announce, then say \emph{ṣarraḥa} again
\end{itemize}

\begin{tcolorbox}[breakable,colback=teal!4,colframe=teal!35!black,title={Before you move on}]
What does \ar{صرح} mean? (\textquotedblleft{}To declare, to state\textquotedblright{}.) Say it once more.
\end{tcolorbox}
